\subsection{A frozen finite predecessor which new ordinary primes cannot repair}

\begin{result}{E.149}{Proved}{The finite ordinary extension gate can genuinely fail}
There are admissible level-one data of E.148
with lower homomorphism equal to the actual
point-\(5/16\) section in
\[
                 Q_1=F/P_2\simeq\mathbb F_3(1)^2
\]
for which no next-level ordinary extension
exists, even if arbitrarily many additional
primes of any congruence class are allowed.

One may take the old finite support to be
\(\{2,5,11,q\}\), where
\[
                    q\equiv2\pmod9,
\]
so all old support primes have
\(v_3(p^2-1)=1\).
At two, five, and eleven the point target
is the original \(5/16\); at \(q\) it is
\(1-q^m\), with any specified positive
integer \(m\) divisible by three.
The point at three may be fixed to \(5/16\)
at the abelian level.
The obstruction is already the plus
Frobenius equation at the two old places
five and \(q\).

The same failed finite data can be
realized by an actual ordinary completed
pair of infinite support, by changing
one central Frobenius prescription in
E.125. This proves failure for specific
frozen predecessors; it does not prove
that every possible infinite extension
path must eventually fail.
\end{result}
\begin{proof}
The first relevant quotient is
\[
                  Q_2/B_2\simeq H_3,
 \tag{E.149.1}
\]
the marked exponent-three Heisenberg
group. Indeed \(Q_2=F/P_3\) has
abelian exponent nine and central
commutator of order three; killing
\(B_2=\langle x^3,y^3\rangle\)
gives the usual group of order \(27\).
Its center \(C\simeq\mathbb F_3(2)\)
has trivial arithmetic action.

Let \(s_0\) be the actual point-\(5/16\)
section in \(H_3\). Retain \(\mu_9\),
the actual action, and the cusp-one
path in a finite arithmetic quotient.
The joint image of \(s_0\) has full
geometric \(H_3\) over that quotient:
inertia at five and eleven is trivial
on the arithmetic data and supplies
the two independent Frattini vectors,
as in E.126.
Choose an arithmetic element \(d\)
with cyclotomic value two modulo nine.
Let \(s_1(d)\) denote the actual
cusp-one value above it, and choose
\(0\ne z\in C\).
The joint image contains
\[
                             z\,s_1(d).
 \tag{E.149.2}
\]
Chebotarev supplies a prime \(q\),
outside every prescribed finite set,
whose joint Frobenius is this element,
up to conjugacy.
Choose a prime over it to use the
displayed representative. Thus
\(q\equiv2\pmod9\).
The local reference class at \(q\)
differs from the actual cusp-one
class by the nonzero central
unramified character with Frobenius
value \(z\).

Freeze the actual point \(1-q^m\),
where \(3\mid m\).
Its \(H_3\)-image is the cusp-one
class: the local cusp twist has
peripheral powers \(y^{m\kappa(q)}\),
which vanish in the exponent-three
group. The identification with the
actual point is the boundary model
and rigidity in A.31.
Its lower \(Q_1\)-class is zero.
The reference point has the same
lower local class, because it is
a unit point at \(q\) and every
\(q\)-adic unit is a cube.
Hence the prescribed finite lower
homomorphism really does match this
new frozen ordinary point.
It also matches all the original
old point targets.
At every other prime outside the
old set, its abelian class is that
of the good-reduction rational
point \(5/16\).
Thus these are admissible data in
the precise sense of E.148.

Put \(S=\{2,3,5,11,q,\infty\}\).
Every prime adjoined to
\(\{3,\infty\}\) here is minus one
modulo three. Finite purity, or
E.142 with even weight, therefore
gives
\[
 H^1(G_{\mathbb Q,S},C)
   =H^1(G_{\mathbb Q,\{3,\infty\}},\mathbb F_3)
   =\mathbb F_3\cdot\bar\eta_4.
 \tag{E.149.3}
\]
In particular adjoining \(q\)
creates no new constant-\(\mathbb F_3\)
character: its tame inertia coinvariant
is killed by the unit \(q-1\).
The remaining global line is the
cyclotomic cubic character.
All \(H_3\)-lifts of the fixed
literal lower homomorphism over
\(G_{\mathbb Q,S}\) differ from \(s_0\)
by this line.

At five the local abelian stabilizer
is zero:
\[
 H^0(\mathbb Q_5,Q_1)=0,
 \]
because Frobenius acts by minus one.
The central fibre action is consequently
free. Since
\[
                       \bar\eta_4(F_5)=1,
 \]
matching the unchanged point target
at five forces the global central
twist to be zero.
At \(q\) the local lower stabilizer
is zero for the same reason.
The nonzero central discrepancy \(z\)
in (E.149.2) therefore cannot be
absorbed by a geometric conjugacy.
The unique globally pinned class
fails the point match there.
In the notation of E.148, every
allowed aligning frame gives
\(\Delta_5^+=0\) and
\(\Delta_q^+=z\ne0\).
Thus the exact plus equation fails.

Suppose one now allowed any further
ordinary support primes and found
a \(Q_2\)-extension retaining this
literal lower homomorphism and all
the frozen point targets.
At each new prime its lower
abelian inertia is zero, so its
ordinary valuation is a multiple
of three. Its inertia in the
quotient \(H_3\) is then trivial.
This includes the infinity
peripheral, by E.144.
The projected \(H_3\)-section is
therefore unramified outside the
old \(S\) and retains both old
point matches, contradicting the
preceding calculation.
This is an obstruction to the
actual finite extension problem,
not merely to an integral
correction or to one choice of
denominator-clearing exponents.

There is also a one-place variant.
Choose the arithmetic element in
(E.149.2) with cyclotomic value
minus one modulo nine instead.
Then
\[
                  \bar\eta_4(F_q)
                        =(q^4-1)/3=0\pmod3.
 \]
Every global central twist in
(E.149.3) is already locally zero
at \(q\), so the nonzero discrepancy
there is by itself an obstruction.
This variant has a nonminimal
three-adic denominator; the
two-place construction above does
not need that feature.

Finally we verify the assertion
about a genuine completed pair.
Run the ordinary E.125 recursion,
but at one fresh \(B\)-support
prime retain the marked abelian
Frobenius coordinates zero modulo
nine and replace its marked
cusp-one central coordinate by
that coordinate plus one modulo
nine. Keep every radical, congruence,
and local-root condition unchanged.
The point-reference finite image
contains the full center.
The central change does not alter
the abelian ninth-root intersection
with the current cyclic radical
extension, so the E.128 compatibility
argument still supplies the prime.
It remains two modulo nine.
The ordinary convergence proof of
E.125 is unchanged.

The integral normalization
\(h\equiv h_{5/16}\pmod9\)
also remains valid: it uses the
nine-divisibility of every fresh
inertia exponent and the old
two/five/eleven comparisons,
not the new central Frobenius
value.
At five the normalized class
matches the reference point and
pins the integral global
three-primary twist to zero.
At the modified prime the exact
boundary discrepancy is one
modulo three; its local group
has order three.
Thus the same two old support
places witness failure.
The ordinary pair has at that
prime the factorial exponent
specified by E.125, which is
divisible by three, and its
full local pro-three point class
is the model used above.
All global abelian reductions
modulo three are those of
\(5/16\).
This realizes the failed finite
predecessor on an ordinary
infinite-support pair, without
asserting failure of any stronger
unrestricted nilpotent lifting
condition not established here.

The exact finite certificate
\path{code/check_failed_predecessor.py}
enumerates all \(27\) geometric
conjugators for each of the three
global central twists.
Its old-five/new-\(q\) matching
counts are respectively
\((3,0),(0,3),(0,0)\), so no twist
matches both.
It also checks the nonminimal
one-place variant and the distinct
central twisted-conjugacy classes.
The saved output is
\path{code/check_failed_predecessor.out}.
This finite certificate checks
the normal forms and conjugacy;
the arithmetic prime existence
and completed-pair realization
are the separate arguments above.
\end{proof}
